%==============================================================================
% Sjabloon poster bachproef
%==============================================================================
% Gebaseerd op document class `a0poster' door Gerlinde Kettl en Matthias Weiser
% Aangepast voor gebruik aan HOGENT door Jens Buysse en Bert Van Vreckem

\documentclass[a0,portrait]{hogent-poster}

% Info over de opleiding
\course{Bachelorproef}
\studyprogramme{toegepaste informatica}
\academicyear{2025-2026}
\institution{Hogeschool Gent, Valentin Vaerwyckweg 1, 9000 Gent}

% Info over de bachelorproef
\title{Cybersecurity-automatisatie in CI/CD-pipelines voor CRA-compliance}
%\subtitle{Ondertitel (eventueel)}
\author{Ilan Engels}
\email{ilan.engels@student.hogent.be}
\supervisor{Pieter-Jan Maenhaut}
\cosupervisor{Moreno Robyn (Excentis)}

% Indien ingevuld, wordt deze informatie toegevoegd aan het einde van de
% abstract. Zet in commentaar als je dit niet wilt.
\specialisation{Systeem- en Netwerkbeheer}
\keywords{CRA, Cybersecurity, CI/CD, Jenkins, Pipeline, DevSecOps, Compliance, SBOM, SCA, SAST, DAST}
\projectrepo{https://github.com/EngelsIlan/BAP}

\begin{document}

\maketitle

\begin{abstract}
  Deze bachelorproef onderzoekt hoe CRA-compliance geautomatiseerd kan worden binnen CI/CD-pipelines. De focus ligt op een Jenkins-omgeving bij Excentis, waar momenteel nog geen securitycontroles worden uitgevoerd. Er werd een Proof of Concept ontwikkeld met geautomatiseerde securitystappen voor SBOM-generatie, dependency scanning, SAST, DAST en secrets scanning. De pipeline gebruikt Syft, OWASP Dependency-Check, SonarQube, Grype, TruffleHog en OWASP ZAP. De resultaten tonen aan dat CRA-relevante securitycontroles technisch kunnen worden geïntegreerd in een bestaande CI/CD-pipeline zonder het ontwikkelproces fundamenteel te verstoren.
\end{abstract}

\begin{multicols}{2} % This is how many columns your poster will be broken into, a portrait poster is generally split into 2 columns

\section{Introductie}

De toenemende digitalisering van softwareontwikkeling brengt niet alleen efficiëntiewinst met zich mee, maar verhoogt ook de nood aan structurele beveiliging. Tegelijk legt de Europese Cyber Resilience Act (CRA) strengere eisen op aan softwareproducenten om digitale producten aantoonbaar veilig en beheersbaar te maken voor ze op de markt worden gebracht. Voor organisaties die software ontwikkelen, betekent dit dat beveiliging niet langer een afzonderlijke eindcontrole kan zijn, maar geïntegreerd moet worden in het volledige ontwikkelproces, en best zo vroeg als mogelijk.

In deze bachelorproef wordt onderzocht hoe CRA-compliance geautomatiseerd kan worden binnen CI/CD-pipelines. De focus ligt op een Jenkins-omgeving bij Excentis, waar momenteel nog geen securitycontroles geïmplementeerd zijn. Dit vormt een relevante context, omdat handmatige controles moeilijk schaalbaar zijn en vaak tijdrovend en kostelijk zijn.

Om dit probleem aan te pakken werd een Proof of Concept ontwikkeld waarin verschillende securitycontroles worden geïntegreerd in een pipeline. Daarbij werd onder meer gebruikgemaakt van SBOM-generatie, dependency scanning, SCA, SAST, DAST en secrets scanning. Voor de implementatie werden Syft, OWASP Dependency-Check, SonarQube, Grype, TruffleHog en OWASP ZAP ingezet.

Het doel van deze bachelorproef is niet enkel om te tonen dat dergelijke controles technisch mogelijk zijn, maar ook om te onderzoeken in hoeverre deze controles met de huidige tools binnen een bestaande CI/CD-pipeline kunnen worden geautomatiseerd. Op die manier wil dit onderzoek een concrete bijdrage leveren aan het ondersteunen van veilige softwareontwikkeling en het voorbereiden op toekomstige CRA-verplichtingen.

\section{Experimenten}

Om te onderzoeken in hoeverre CRA-relevante securitycontroles geautomatiseerd kunnen worden binnen een bestaande CI/CD-pipeline werd een Proof of Concept opgesteld en getest in een Jenkins-omgeving. De experimenten hadden als doel na te gaan of verschillende beveiligingsstappen, zoals SBOM-generatie, dependency scanning, SCA, SAST, DAST en secrets scanning op een stabiele en bruikbare manier konden worden geïntegreerd.

De pipeline werd uitgeprobeerd op representatieve testprojecten om de werking van de afzonderlijke securitytools en hun onderlinge samenhang te evalueren. Daarbij werd onder meer gekeken naar de uitvoerbaarheid van de scans, de kwaliteit van de rapporten, de impact op de buildflow en de mate waarin kwetsbaarheden automatisch konden worden opgespoord en gerapporteerd.

Tijdens de experimenten werd vastgesteld dat de geselecteerde tools technisch combineerbaar zijn binnen één enkele pipeline, al bracht dit ook enkele uitdagingen met zich mee, zoals uitvoertijd, toolintegratie en false positives. De experimenten tonen aan dat geautomatiseerde CRA-controles in een CI/CD-omgeving technisch haalbaar zijn, maar op voorwaarde dat de pipeline correct wordt opgebouwd en afgestemd op de gebruikte technologieën.

\section{Sectie met figuur}

De {\LaTeX} figure-omgeving bepaalt zelf waar een afbeelding komt en dat is meestal niet op de plek in de tekst waar de figure-omgeving gedefinieerd wordt. Als je wilt forceren dat afbeeldingen toch in de flow van de tekst blijven, dan kan je dat zoals hieronder:

\begin{center}
  \captionsetup{type=figure}
  \includegraphics[width=1.0\linewidth]{grail}
  \captionof{figure}{He hasn't got shit all over him. The nose? Where'd you get the coconuts? What do you mean? We shall say `Ni' again to you, if you do not appease us}
\end{center}

Let er wel op dat dit tot problemen met bladschikking kan leiden.

\section{Conclusies}

Don't underestimate the Force. Oh God, my uncle. How am I ever gonna explain this? I suggest you try it again, Luke. This time, let go your conscious self and act on instinct. Don't be too proud of this technological terror you've constructed. The ability to destroy a planet is insignificant next to the power of the Force.

\section{Toekomstig onderzoek}

I care. So, what do you think of her, Han? No! Alderaan is peaceful. We have no weapons. You can't possibly… I have traced the Rebel spies to her. Now she is my only link to finding their secret base.

Kid, I've flown from one side of this galaxy to the other. I've seen a lot of strange stuff, but I've never seen anything to make me believe there's one all-powerful Force controlling everything. There's no mystical energy field that controls my destiny. It's all a lot of simple tricks and nonsense. You are a part of the Rebel Alliance and a traitor! Take her away! 

\end{multicols}
\end{document}